\documentclass[10pt,a4paper]{amsart}
\usepackage{amsmath,amssymb,amsthm,graphicx}
\usepackage{verbatim}
\newenvironment{reference}{\comment}{\endcomment}
%\newenvironment{reference}{}{}
\newenvironment{slogan}{\comment}{\endcomment}
\newenvironment{history}{\comment}{\endcomment}

% For commutative diagrams we use Xy-pic
\usepackage[all]{xy}

% We use 2cell for 2-commutative diagrams.
\xyoption{2cell}
% Theorem environments.
%
\theoremstyle{plain}
\newtheorem{theorem}[subsection]{Theorem}
\newtheorem{proposition}[subsection]{Proposition}
\newtheorem{lemma}[subsection]{Lemma}

\theoremstyle{definition}
\newtheorem{definition}[subsection]{Definition}
\newtheorem{example}[subsection]{Example}
\newtheorem{exercise}[subsection]{Exercise}
\newtheorem{situation}[subsection]{Situation}

\theoremstyle{remark}
\newtheorem{remark}[subsection]{Remark}
\newtheorem{remarks}[subsection]{Remarks}

\numberwithin{equation}{subsection}

% Macros
%
\def\lim{\mathop{\mathrm{lim}}\nolimits}
\def\colim{\mathop{\mathrm{colim}}\nolimits}
\def\Spec{\mathop{\mathrm{Spec}}}
\def\Hom{\mathop{\mathrm{Hom}}\nolimits}
\def\Ext{\mathop{\mathrm{Ext}}\nolimits}
\def\SheafHom{\mathop{\mathcal{H}\!\mathit{om}}\nolimits}
\def\SheafExt{\mathop{\mathcal{E}\!\mathit{xt}}\nolimits}
\def\Sch{\mathit{Sch}}
\def\Mor{\mathop{\mathrm{Mor}}\nolimits}
\def\Ob{\mathop{\mathrm{Ob}}\nolimits}
\def\Sh{\mathop{\mathit{Sh}}\nolimits}
\def\NL{\mathop{N\!L}\nolimits}
\def\CH{\mathop{\mathrm{CH}}\nolimits}
\def\proetale{{pro\text{-}\acute{e}tale}}
\def\etale{{\acute{e}tale}}
\def\QCoh{\mathit{QCoh}}
\def\Ker{\mathop{\mathrm{Ker}}}
\def\Im{\mathop{\mathrm{Im}}}
\def\Coker{\mathop{\mathrm{Coker}}}
\def\Coim{\mathop{\mathrm{Coim}}}

% Boxtimes
%
\DeclareMathSymbol{\boxtimes}{\mathbin}{AMSa}{"02}

%
% Macros for moduli stacks/spaces
%
\def\QCohstack{\mathcal{QC}\!\mathit{oh}}
\def\Cohstack{\mathcal{C}\!\mathit{oh}}
\def\Spacesstack{\mathcal{S}\!\mathit{paces}}
\def\Quotfunctor{\mathrm{Quot}}
\def\Hilbfunctor{\mathrm{Hilb}}
\def\Curvesstack{\mathcal{C}\!\mathit{urves}}
\def\Polarizedstack{\mathcal{P}\!\mathit{olarized}}
\def\Complexesstack{\mathcal{C}\!\mathit{omplexes}}
% \Pic is the operator that assigns to X its picard group, usage \Pic(X)
% \Picardstack_{X/B} denotes the Picard stack of X over B
% \Picardfunctor_{X/B} denotes the Picard functor of X over B
\def\Pic{\mathop{\mathrm{Pic}}\nolimits}
\def\Picardstack{\mathcal{P}\!\mathit{ic}}
\def\Picardfunctor{\mathrm{Pic}}
\def\Deformationcategory{\mathcal{D}\!\mathit{ef}}
\usepackage{cleveref}
\crefname{originalreference}{original reference}{original references}
\Crefname{originalreference}{Original reference}{Original references}
\makeatletter
\newlabel{sites-section-sheaves-algebraic-structures}{{00YR}{source}{}{sourcepointer-1}{}}
\newlabel{sites-section-sheaves-algebraic-structures@cref}{{[originalreference][0][]00YR}{[1][0][]source}{}{}{}}
\makeatother
\crefname{originalreference}{original reference}{original references}
\Crefname{originalreference}{Original reference}{Original references}
\makeatletter
\newlabel{sites-definition-topos}{{00XA}{source}{}{sourcepointer-2}{}}
\newlabel{sites-definition-topos@cref}{{[originalreference][0][]00XA}{[1][0][]source}{}{}{}}
\makeatother
\title{Exactness of the inverse image on abelian sheaves and of the pullback along a flat morphism of ringed topoi}
\author{The Stacks Project Authors}
\date{stacks-project:sites-modules.tex:local-snapshot-20261012; GNU FDL 1.2 or later}
\begin{document}
\maketitle
\noindent\textit{Editorial note.} All delivered mathematical text is the author's own, copied verbatim from the pinned Stacks Project chapter file sites-modules.tex; no statement or proof step was rewritten, reordered or removed.\par\medskip
\noindent\textit{Editorial note.} The counted span cites two external items (a sites section and the definition of a morphism of topoi), recorded in assembly.external\_references with canonical tags as navigation pointers.\par\medskip
\noindent\textit{Allowed standard prerequisites (editorial).}\begin{itemize}
\item Definition of a morphism of topoi and exactness of the inverse image on sheaves of sets. Stacks Project, Sites, Definition sites-definition-topos (tag 00X4). Supplies the exactness of the inverse image on sets from which the abelian-sheaf statement follows.
\item Flat morphisms of ringed topoi and the definition of pullback of modules. Stacks Project, Modules on Sites, section on flatness (definition of a flat morphism of ringed topoi). Turns exactness of the pullback into the flatness hypothesis.
\end{itemize}\medskip
\ifdefined\hypertarget\hypertarget{sourcepointer-1}{}\fi
\noindent\textit{Source reference.} Original reference 00YR: sources/stacks-chow/tags.txt, source lines 1265--1265. Cross-chapter standard prerequisite cited by the author inside the counted span; it is used as a pointer and its own proof is not part of this item.\par\medskip
\ifdefined\hypertarget\hypertarget{sourcepointer-2}{}\fi
\noindent\textit{Source reference.} Original reference 00XA: sources/stacks-chow/tags.txt, source lines 1212--1212. Cross-chapter standard prerequisite cited by the author inside the counted span; it is used as a pointer and its own proof is not part of this item.\par\medskip
\begin{lemma}
\label{lemma-flat-pullback-exact}
Let $f : \Sh(\mathcal{C}) \to \Sh(\mathcal{C}')$
be a morphism of ringed topoi. Then
$$
f^{-1} : \textit{Ab}(\mathcal{C}') \longrightarrow \textit{Ab}(\mathcal{C}),
\quad
\mathcal{F} \longmapsto f^{-1}\mathcal{F}
$$
is exact. If
$(f, f^\sharp) :
(\Sh(\mathcal{C}), \mathcal{O})
\to
(\Sh(\mathcal{C}'), \mathcal{O}')$
is a flat morphism of ringed topoi then
$$
f^* : \textit{Mod}(\mathcal{O}') \longrightarrow \textit{Mod}(\mathcal{O}),
\quad
\mathcal{F} \longmapsto f^*\mathcal{F}
$$
is exact.
\end{lemma}

\begin{proof}
Given an abelian sheaf $\mathcal{G}$ on $\mathcal{C}'$
the underlying sheaf of sets of $f^{-1}\mathcal{G}$ is the same
as $f^{-1}$ of the underlying sheaf of sets of $\mathcal{G}$, see
Sites, Section \ref{sites-section-sheaves-algebraic-structures}.
Hence the exactness of $f^{-1}$ for sheaves of sets (required in the
definition of a morphism of topoi, see
Sites, Definition \ref{sites-definition-topos})
implies the exactness of $f^{-1}$ as a functor on abelian sheaves.

\medskip\noindent
To see the statement on modules recall that $f^*\mathcal{F}$ is defined
as the tensor product
$f^{-1}\mathcal{F} \otimes_{f^{-1}\mathcal{O}', f^\sharp} \mathcal{O}$.
Hence $f^*$ is a composition of functors both of which are exact.
\end{proof}

\end{document}
